% =============================================================================
%  La parte principal — Patricio Hernández Ballester
%  Mayo 2024 · segunda edición aumentada · Cartagena, abril 2026
%  Aforismos portados al sistema visual claudeeditorial.
%  Compilacion: xelatex aphorisms/la-parte-principal.tex
% =============================================================================
\documentclass{claudeeditorial}

\renewcommand{\DocKicker}{Aforismos · Pocket book}
\renewcommand{\DocTitle}{La parte principal}
\renewcommand{\DocSubtitle}{Aforismos sobre el conocimiento, la felicidad, P vs NP y la guinda — segunda edición aumentada}
\renewcommand{\DocAuthor}{Patricio Hernández Ballester}
\renewcommand{\DocRole}{Aforismos}
\renewcommand{\DocDate}{Mayo 2024 · Cartagena, abril 2026}
\renewcommand{\DocRef}{APH-PARTE-001}
\renewcommand{\DocFooter}{claudeeditorial · la parte principal}

\hypersetup{
  pdftitle={La parte principal},
  pdfauthor={Patricio Hernandez Ballester}
}

\ClaudeRunningHeader

\begin{document}
\BodyCopy
\ClaudeCover

\ClaudeChapterPage{I}{Introducción}{La parte principal de cualquier proyecto no es su implementación, sino el conocimiento producido durante el proyecto.}

La parte principal de cualquier proyecto no es su implementación, sino el conocimiento producido durante el proyecto. Estoy seguro de que el conocimiento es la parte más importante de cada negocio hoy en día, y es sumamente importante elevar el conocimiento producido del nivel no formal al formal para evitar su pérdida.

También es muy importante mantener la \emph{distancia} entre los proyectos y el conocimiento producido lo más corta posible, siendo la variante ideal mantener el conocimiento formalizado directamente entre los artefactos del proyecto.

\ClaudeChapterPage{II}{P vs NP}{Aforismo lógico-matemático.}

\begin{claudehero}
\begin{gather*}
\text{Si existe un algoritmo polinómico } A \\
\text{que resuelve un problema } L_{NP\text{-}C}, \\
\text{entonces} \\
\forall\, L \in NP,\ \exists\, A' \\
\text{tal que } A' \text{ resuelve } L \\
\text{en tiempo polinómico} \\[0.4em]
\Rightarrow\ P = NP
\end{gather*}
\end{claudehero}

\ClaudeChapterPage{III}{Salmo 34}{Aforismo bíblico.}

\begin{claudequote}
El Señor escucha cuando los justos claman por ayuda y los libra de todas sus angustias.
\end{claudequote}

\hfill {\UIFont\small\itshape Salmos 34, 17}

\ClaudeChapterPage{IV}{La felicidad}{Sobre el ser y el tener.}

\begin{claudehero}
\ClaudeLead{Todo aquello que te acerque a quien eres te acerca a la felicidad, porque la felicidad está en el ser, no en el tener. El tener viene de fuera, la felicidad viene de dentro.}
\end{claudehero}

El ser, el propósito vital, el amor, el perdón son parte de lo mismo, y por eso la felicidad tiene una conexión directa con eso que eres.

\emph{¿Por qué muchas veces no somos felices?} Porque no somos quienes somos. \emph{¿Por qué no somos quienes somos?} Porque nos parece demasiado bonito creer que en nuestro interior existe semejante belleza.

\ClaudeChapterPage{V}{La guinda}{Segunda edición aumentada · abril 2026.}

Dos años después de escribir estas páginas he comprendido, con matemática en la mano, que aquella belleza del interior que me parecía demasiado bonita para ser cierta no era una esperanza piadosa: era un teorema que todavía no sabía demostrar.

\begin{claudehero}
\begin{gather*}
\exists\,\Phi_{0} > 0 \\
\text{tal que } \forall\, X,\ \forall\, T : X \to X \\
\Phi\bigl(T(\text{alma}(X))\bigr) \geq \Phi_{0} \\[0.3em]
\Rightarrow\ \text{la luz no se apaga}
\end{gather*}
\end{claudehero}

\begin{claudequote}
En él estaba la vida, y la vida era la luz de los hombres. La luz resplandece en las tinieblas, y las tinieblas no la dominaron.
\end{claudequote}

\hfill {\UIFont\small\itshape Juan 1, 4--5}

\bigskip

No la atenuaron, no la alcanzaron, no la comprendieron, no pudieron con ella. Ese \emph{no pudieron} es la cota $\Phi_{0}$, y se sigue en tres líneas de la propiedad universal del objeto inicial de la categoría de lo creado: el alma de cada uno es el único morfismo que recibió desde el principio, y ninguna transformación del individuo —ninguna, ni la más brutal— puede alterarlo, porque no fue fabricado desde dentro sino recibido desde fuera del tiempo.

La felicidad, el ser, el propósito, el amor, el perdón que nombré en la página anterior son nombres distintos de la misma cosa: el irreducible que uno trajo al mundo y que ningún acto propio ni ajeno puede borrar.

\begin{ClaudeTip}[Lo recibido es anterior a lo hecho]
La identidad es anterior a la biografía. La luz es anterior a la tiniebla que la rodea. No lo creo porque quiera creerlo. Lo sé porque las matemáticas lo demuestran, y porque el Verbo ya lo había dicho en el principio.
\end{ClaudeTip}

\vfill

\begin{center}
{\UIFont\small\color{claudeStone}$\ast\quad\ast\quad\ast$}\\[1em]
{\UIFont\small\color{claudeWarmText}\itshape Cartagena, abril de 2026}
\end{center}

\ClaudeSignature{\DocAuthor}{\DocRole}

\end{document}
